\section{Introduction}\label{sec:introduction}

For bounded sets $A,B$ in a finite-dimensional real vector space, with
$B$ having nonempty interior, let $N(A,B)$ denote the least number of
translates of $B$ that cover $A$. The translation vectors may lie anywhere
in the ambient space. For an origin-symmetric convex body
$K\subset\R^n$, its polar is
\[
 K^\circ=\{y\in\R^n:\langle x,y\rangle\le1
                    \text{ for every }x\in K\}.
\]
Here a convex body is compact and has nonempty interior. All logarithms
are natural.

The duality conjecture for metric entropy, originating in Pietsch's
operator-theoretic work in 1972, asks whether there are absolute constants
$a,b\ge1$ such that
\begin{equation}\label{eq:duality-conjecture}
 \frac1b\log N(L^\circ,aK^\circ)
 \le \log N(K,L)
 \le b\log N(L^\circ,a^{-1}K^\circ)
\end{equation}
for every dimension and every pair of origin-symmetric convex bodies
$K,L$; see \cite[Conjecture~1]{AMSTJ2004}.
It asks whether passing to the polar pair preserves covering complexity
up to universal changes of scale and multiplicative constants.

We disprove the full conjecture.

\begin{theorem}\label{thm:main}
For every $a,b\ge1$, there are a positive integer $n$ and an
origin-symmetric convex body $K\subset\R^n$ such that, with
$L=[-1,1]^n$,
\[
 \log N(K,L)>b\log N(L^\circ,a^{-1}K^\circ).
\]
\end{theorem}

Thus the upper bound in \eqref{eq:duality-conjecture} fails for every
proposed pair of absolute constants. In particular, the two-sided
duality conjecture has a negative answer. The bodies in
Theorem~\ref{thm:main} have nonempty interior in the full ambient
dimension. Their dimensions grow with the parameters; no assertion
about a fixed dimension is intended.

\subsection{Related work}

The question developed from the duality problem for entropy numbers of
operators; Artstein, Milman and Szarek give the geometric and operator
formulations in \cite[pp.~1314--1315]{AMS2004} and attribute the original
problem to Pietsch's 1972 work. Early progress already showed that
polarity controls covering numbers in several regimes. K\"onig and Milman
proved a comparison of the $n$th roots of the two covering numbers
\cite[Theorem~1]{KonigMilman1987}; on the logarithmic scale this allows
an additive error proportional to $n$. Bourgain, Pajor, Szarek and
Tomczak-Jaegermann established entropy-number comparisons under
geometric and sequence-regularity assumptions, as well as estimates
with losses depending on the rank \cite{BPSTJ1989}.

Artstein, Milman and Szarek proved the dimension-free conjecture when
one body is a Euclidean ball, and hence when either body is an ellipsoid
\cite[Theorem~1]{AMS2004}. Artstein, Milman, Szarek and
Tomczak-Jaegermann proved universal duality for \emph{convexified packing}
\cite[Theorem~2]{AMSTJ2004}. This quantity uses an ordered sequence of
points: the interior of the translate centered at each new point must
avoid the convex hull of its predecessors. That requirement is stronger
than pairwise separation and differs from the ordinary covering number
in \eqref{eq:duality-conjecture}. They also announced ordinary entropy
duality for spaces whose K-convexity constant is bounded, with the
comparison constants depending on that bound
\cite[Theorem~5]{AMSTJ2004}; see also the later derivation in
\cite[Section~2]{Milman2007}. A bound for each individual
finite-dimensional space does not provide one pair of constants for
all spaces and dimensions.

For arbitrary symmetric convex bodies, Emanuel Milman obtained a
comparison with logarithmic losses in both the scale and the entropy
factor \cite[Corollary~1.2]{Milman2007}. Connections with learning theory
provide another perspective: Hu, Peale and Reingold study dual covering
numbers of function classes, obtaining a comparison with an explicit
prefactor depending on the approximation scale and the uniform bounds
of the function classes \cite[Theorem~11]{HuPealeReingold2022}.
These dimension- or scale-dependent estimates are compatible with the
failure of a universal pair $a,b$.

The column approximation in this paper also belongs to the broader
study of entropy of convex hulls. For example, Mendelson develops
finite-set convex-hull estimates in $L_2(\mu)$ using Maurey's
approximation method
\cite[Section~2.1]{Mendelson2001}. Here the approximation is uniform over
all row coordinates and uses the special partition structure: labels
replace the original centers in the finite encoding. The finite-field
side uses the polynomial zero estimate associated with Schwartz and
Zippel \cite{Schwartz1980,Zippel1979}, together with diagonal recovery
of symmetric multilinear forms in sufficiently large characteristic
\cite{FerreroMicali1979,DrapalVojtechovsky2009}. We prove the precise
elementary forms of these tools and the additional compression and
separation arguments below.

Theorem~\ref{thm:main} also separates ordinary covering entropy from
convexified-packing entropy, even after a fixed change of scale;
Corollary~\ref{cor:convexified-separation} states and proves this
consequence after the construction.

\subsection{The construction and its main mechanism}

The proof starts with a finite real matrix whose entries lie in $[0,1]$.
Its rows are pairwise separated by one in the sup norm, while the
absolutely convex hull of its columns admits a small uniform
approximation list. A direct support-function calculation converts
these two properties into large primal entropy and small polar entropy.

The approximation statement is the reusable part of the argument.
Suppose a finite set carries $u$ partitions, each with at most $q$ classes.
For a subset $T$ of the partition indices, join two points whenever they
belong to the same class in one partition indexed by $T$. We consider
the resulting graph distances $d_T$, allowing infinite distance between
components. For a positive integer cutoff $h$, define the profiles
\[
 x\longmapsto
 \max\left\{0,1-\frac{\log(1+d_T(x,v))}{\log(1+h)}\right\},
\]
with value zero at infinite distance.
The columns consist of these profiles over all centers $v$ and over a
family of sets $T$, each omitting only a small fraction of the partitions.

Proposition~\ref{prop:compression} gives a uniform approximation of
every nonnegative combination of these columns whose coefficients sum
to at most one. For a given combination,
averaging finds one partition present in all but a small amount of
coefficient mass. At each distance level, a greedy selection of a bounded
number of pivots captures all but a small residual mass at every point.
Each pivot can then be replaced by a fixed representative of its class
in the common partition. This replacement requires storing a label,
rather than a point of the original set.

The replacement enlarges the relevant distance radii from $j$ to at
most $3j+2$. The logarithmic profile bounds the error caused by this enlargement
by $\log3/\log(1+h)$, uniformly over all columns and all graph
components. After grouping weights by pivot slot and by $T$, the
approximation count depends on $q$ through a fixed power. It does not
depend on the number of possible column centers through an additional
factor. Both the choice of the common partition and all grouped
weights are included in the count. Splitting a signed coefficient vector
into its positive and negative parts then gives an approximation list
for the absolutely convex column hull, with twice the full error of the nonnegative
approximation and the square of the list size.

To obtain separated rows, we use symmetric $h$-linear forms on
$\F_p^r$, where $p$ is a sufficiently large prime and $r$ is a positive
integer. A probabilistic choice of directions ensures that, for
every nonzero form, one can delete a small fraction of the directions
so that no evaluation on a remaining $h$-tuple vanishes, including
tuples with repetitions.
Contraction with one direction supplies a partition. A short path in the
associated graph would express a nonzero form as a sum of forms killed
by the directions along the path, contradicting the nonvanishing
property.

There are $p^{D_h}$ rows but at most $p^{D_{h-1}}$ labels in each
partition, where $D_j=\binom{r+j-1}{j}$ for $j\ge0$. The ratio
\[
 \frac{D_{h-1}}{D_h}=\frac{h}{r+h-1}
\]
tends to zero as $r$ grows with $h$ fixed. All remaining encoding costs
are fixed before the prime $p$ is chosen. Taking $p$ sufficiently large
therefore makes those costs negligible. This order of choices is what
allows the compression result and the separation construction to
produce a failure of dimension-free entropy duality.

Section~\ref{sec:convex-bodies} first proves the matrix-to-body conversion,
which serves as the construction target. Section~\ref{sec:compression}
then establishes uniform compression for an arbitrary admissible family
of partitions. Section~\ref{sec:partitions} constructs the finite-field
partitions and proves row separation. Finally, Section~\ref{sec:parameters}
chooses the parameters and applies the conversion to defeat any proposed
universal constants, then derives the convexified-packing consequence.
All arguments needed for the construction are included.

\subsection{Notation}

For a finite set $S$, write
\[
 B_\infty^S=[-1,1]^S,\qquad
 B_1^S=\left\{\lambda\in\R^S:
                   \sum_{s\in S}|\lambda_s|\le1\right\}.
\]
The sup norm on functions on $S$ is denoted by $\|\cdot\|_{\infty,S}$.
For a finite collection of vectors, $\absconv$ denotes the convex hull
of the vectors and their negatives. A uniform $\varepsilon$-approximation
list may have centers outside the class being approximated; actual
centers in the coefficient ball will be chosen in
Lemma~\ref{lem:matrix-to-bodies}. We write $[u]=\{1,\ldots,u\}$.
